\chapter{Results and Discussion}

% =====================================================================
% Ordenado por eje de medida, no por configuración: así cada sección es ya
% una comparación entre las cuatro alternativas. Ninguna definición de métrica
% en este capítulo; todas están en el Capítulo 5.
% =====================================================================

\section{Delivery Performance}
\emph{Presentar la proporción de datos entregados en las cuatro configuraciones y
el reparto de la carga entre interfaces. Es la sección que establece que ninguna
alternativa gana rendimiento a costa de perder información, condición previa para
que las comparaciones siguientes sean legítimas.}

\section{Latency and Information Freshness}
\emph{Presentar el retardo extremo a extremo y la frescura del dato en recepción,
con especial atención a la proporción de datos que llegan dentro de plazo. Aquí
debe verse el efecto diferencial de la extensión por plazos frente a la política
predictiva sin ella.}

\section{Buffer Occupancy and Data Loss}
\emph{Presentar la evolución de la ocupación de memoria y los descartes
producidos. Es la evidencia directa de que la condición de supervivencia
formulada en el Capítulo 3 se cumple en ejecución, y donde la referencia de
acceso inalámbrico exclusivo debe mostrar su límite.}

\section{Energy Consumption}
\emph{Presentar el consumo por dominio energético y la potencia media,
comparando las cuatro configuraciones. Cuantificar el ahorro que aporta la
retención predictiva frente al uso permanente del acceso celular.}

\section{Energy--Freshness Trade-off Across Configurations}
\emph{Sección donde aterriza la tesis. Cruzar coste energético y frescura de la
información para situar cada configuración en el espacio de compromiso y mostrar
qué posición ocupa la propuesta completa. Debe responder a por qué la extensión
por plazos justifica su coste energético adicional.}

\section{Discussion}
\emph{Interpretar los resultados en conjunto: qué comportamientos confirman las
propiedades esperadas de la política, cuáles no se anticiparon, y cómo se
explican a partir del modelo del Capítulo 3 y de las decisiones de realización
del Capítulo 4.}

\section{Threats to Validity}
\emph{Declarar explícitamente los límites de validez de los resultados: perfiles
de predicción derivados del propio escenario, ruta única, modelo energético
basado en especificaciones y no en medidas de hardware, y ausencia de
interferencia externa. Anticipar estas objeciones aquí las desactiva.}
